En este capítulo se muestran ejemplos de los elementos que generan listas automáticas.

\section{Ejemplo de Figuras}
texto de ejemplo describiendo la figura.

\begin{figure}[h]
    \centering
    % Usamos \rule para simular una imagen sin necesitar un fichero externo
    \rule{10cm}{5cm} 
    \caption{Ejemplo de figura simulada (caja negra)}
    \label{fig:ejemplo_figura}
\end{figure}

Como se ve en la \cref{fig:ejemplo_figura}, el sistema funciona.

\section{Ejemplo de Tablas}
texto de ejemplo sobre tablas.

\begin{table}[h]
    \centering
    \begin{tabular}{|c|c|}
        \hline
        \textbf{Dato A} & \textbf{Dato B} \\
        \hline
        1 & 2 \\
        3 & 4 \\
        \hline
    \end{tabular}
    \caption{Tabla de ejemplo para el listado}
    \label{tab:ejemplo_tabla}
\end{table}

\section{Ejemplo de Ecuaciones}
texto de ejemplo matemático.

\begin{equation}
    E = mc^2
    \label{eq:einstein}
\end{equation}

\section{Ejemplo de Algoritmos}
La plantilla usa el entorno \texttt{algorithmN} para algoritmos numerados.

\begin{algorithmN}[alg:mi_algoritmo]{Algoritmo de ejemplo}{Descripción larga del algoritmo para el índice.}
    \SetKwData{Input}{input}
    \SetKwData{Output}{output}
    \Input{Un dato de entrada}
    \Output{Un resultado}
    \BlankLine
    \If{condición verdadera}{
        ejecutar instrucción\;
    }
    \Else{
        ejecutar otra cosa\;
    }
    \Return{resultado}\;
\end{algorithmN}

\section{Ejemplo de Código}
La plantilla usa comandos específicos como \texttt{\textbackslash CCode}. Requiere el fichero en la carpeta \texttt{codes}.

% Sintaxis: \CCode[label]{Caption Corto}{Caption Largo}{fichero}{linea_inicio}{linea_fin}{numero_primera_linea}
\CCode[cod:hola_mundo]{Ejemplo de código C}{Este es un programa simple en C que imprime Hola Mundo.}{ejemplo.c}{1}{6}{1}

\section{Citas bibliográficas}
texto de ejemplo citando a \cite{ejemplo2026}.